\subsection{Variadic Keyword Parameters}

A variadic keyword parameter collects keyword arguments that were not bound to ordinary named parameters. The syntax uses two stars in the function definition:

\begin{verbatim}
def func(**kwargs):
    ...
\end{verbatim}

The name \texttt{kwargs} is only a convention. The two stars are the important syntax. Python programmers usually use \texttt{kwargs} because it means ``keyword arguments collected here.''

The collected object is a normal dictionary.

\subsubsection{Argument Packing Mechanics: Collecting Keyword Overflow into \texttt{**kwargs}}

A function can collect arbitrary keyword arguments using \texttt{**kwargs}.

\begin{verbatim}
def show_kwargs(**kwargs):
    print(f"kwargs:       {kwargs}")
    print(f"type(kwargs): {type(kwargs)}")

show_kwargs(name="Homer", score=42, mood="hungry")
\end{verbatim}

Possible output:

\begin{verbatim}
kwargs:       {'name': 'Homer', 'score': 42, 'mood': 'hungry'}
type(kwargs): <class 'dict'>
\end{verbatim}

The call supplied three keyword arguments. The function definition had one variadic keyword parameter. Python packed the keyword names and their argument objects into one dictionary.

The binding shape is:

\begin{verbatim}
show_kwargs(name="Homer", score=42, mood="hungry")

inside function frame:
    kwargs -> {
        "name": "Homer",
        "score": 42,
        "mood": "hungry"
    }
\end{verbatim}

The dictionary is a normal dictionary object.

\begin{verbatim}
def show_kwargs(**kwargs):
    print(f"type(kwargs): {type(kwargs)}")
    print()

    selected_names = [
        "keys",
        "values",
        "items",
        "get",
        "update",
        "__getitem__",
        "__setitem__",
    ]

    for name in selected_names:
        print(f"{name:<14} {name in dir(kwargs)}")

    print()
    print(f"keys:              {list(kwargs.keys())}")
    print(f"items:             {list(kwargs.items())}")
    print(f"kwargs.get('score'): {kwargs.get('score')}")

show_kwargs(name="Homer", score=42, mood="hungry")
\end{verbatim}

Possible output:

\begin{verbatim}
type(kwargs): <class 'dict'>

keys           True
values         True
items          True
get            True
update         True
__getitem__    True
__setitem__    True

keys:              ['name', 'score', 'mood']
items:             [('name', 'Homer'), ('score', 42), ('mood', 'hungry')]
kwargs.get('score'): 42
\end{verbatim}

Unlike \texttt{*args}, which creates a tuple, \texttt{**kwargs} creates a dictionary. The keys are the keyword names. The values are the objects supplied by the caller.

A function can combine ordinary parameters with \texttt{**kwargs}.

\begin{verbatim}
def show_report(title, **options):
    print(f"title:   {title!r}")
    print(f"options: {options}")
    print(f"type(options): {type(options)}")

show_report(
    "Springfield report",
    score=42,
    mood="calm",
    quote="D'oh!",
)
\end{verbatim}

Possible output:

\begin{verbatim}
title:   'Springfield report'
options: {'score': 42, 'mood': 'calm', 'quote': "D'oh!"}
type(options): <class 'dict'>
\end{verbatim}

The keyword \texttt{title} is bound to the ordinary parameter \texttt{title}. The other keyword arguments are collected into \texttt{options}.

The binding shape is:

\begin{verbatim}
title   -> "Springfield report"
options -> {
    "score": 42,
    "mood": "calm",
    "quote": "D'oh!"
}
\end{verbatim}

If there is no keyword overflow, Python creates an empty dictionary.

\begin{verbatim}
def show_report(title, **options):
    print(f"title:        {title!r}")
    print(f"options:      {options}")
    print(f"type(options): {type(options)}")
    print(f"len(options):  {len(options)}")

show_report("Empty report")
\end{verbatim}

Possible output:

\begin{verbatim}
title:        'Empty report'
options:      {}
type(options): <class 'dict'>
len(options):  0
\end{verbatim}

An empty \texttt{**kwargs} value is not \texttt{None}. It is an empty dictionary.

The parameter name does not have to be \texttt{kwargs}.

\begin{verbatim}
def show_settings(**settings):
    print(f"settings:       {settings}")
    print(f"type(settings): {type(settings)}")

show_settings(theme="light", font_size=14, answer=42)
\end{verbatim}

Possible output:

\begin{verbatim}
settings:       {'theme': 'light', 'font_size': 14, 'answer': 42}
type(settings): <class 'dict'>
\end{verbatim}

The two stars make the parameter variadic. The name only chooses the local name inside the function frame.

The standard \texttt{inspect} module shows this parameter kind explicitly.

\begin{verbatim}
import inspect

def show_report(title, **options):
    print(title, options)

sig = inspect.signature(show_report)

print(f"signature: {sig}")
print()

for name, param in sig.parameters.items():
    print(f"{name:<8} kind={param.kind}")
\end{verbatim}

Possible output:

\begin{verbatim}
signature: (title, **options)

title    kind=POSITIONAL_OR_KEYWORD
options  kind=VAR_KEYWORD
\end{verbatim}

The \texttt{VAR\_KEYWORD} kind means that the parameter collects keyword overflow.

A function can have both \texttt{*args} and \texttt{**kwargs}.

\begin{verbatim}
def show_call(*args, **kwargs):
    print(f"args:       {args}")
    print(f"type(args): {type(args)}")
    print()
    print(f"kwargs:       {kwargs}")
    print(f"type(kwargs): {type(kwargs)}")

show_call("Homer", "Bart", score=42, mood="restless")
\end{verbatim}

Possible output:

\begin{verbatim}
args:       ('Homer', 'Bart')
type(args): <class 'tuple'>

kwargs:       {'score': 42, 'mood': 'restless'}
type(kwargs): <class 'dict'>
\end{verbatim}

The positional overflow goes into the tuple. The keyword overflow goes into the dictionary.

The binding shape is:

\begin{verbatim}
args   -> ("Homer", "Bart")
kwargs -> {
    "score": 42,
    "mood": "restless"
}
\end{verbatim}

The order in the function definition matters. A typical flexible function has this shape:

\begin{verbatim}
def func(required, *args, named=default, **kwargs):
    ...
\end{verbatim}

The meaning is:

\begin{itemize}
    \item ordinary parameters receive explicitly declared arguments;
    \item \texttt{*args} receives extra positional arguments;
    \item keyword-only parameters receive explicitly named arguments after \texttt{*args};
    \item \texttt{**kwargs} receives extra keyword arguments.
\end{itemize}

Here is a small inspection example.

\begin{verbatim}
import inspect

def demo(a, b=42, *args, mode="calm", **kwargs):
    print(a, b, args, mode, kwargs)

sig = inspect.signature(demo)

print(f"signature: {sig}")
print()

for name, param in sig.parameters.items():
    print(f"{name:<8} kind={param.kind} default={param.default}")
\end{verbatim}

Possible output:

\begin{verbatim}
signature: (a, b=42, *args, mode='calm', **kwargs)

a        kind=POSITIONAL_OR_KEYWORD default=<class 'inspect._empty'>
b        kind=POSITIONAL_OR_KEYWORD default=42
args     kind=VAR_POSITIONAL default=<class 'inspect._empty'>
mode     kind=KEYWORD_ONLY default=calm
kwargs   kind=VAR_KEYWORD default=<class 'inspect._empty'>
\end{verbatim}

This signature says exactly how the entrance into the function frame is constructed.

\subsubsection{Argument Unpacking Operations: Expanding Mappings Across Function Call Boundaries with \texttt{**}}

The double star has a second use at the call site. In a function definition, \texttt{**kwargs} packs keyword overflow into a dictionary. In a function call, \texttt{**some\_mapping} unpacks a mapping into separate keyword arguments.

These are opposite directions:

\begin{verbatim}
definition side:
    def func(**kwargs):
        pack keyword overflow into one dictionary

call side:
    func(**mapping)
        unpack one mapping into many keyword arguments
\end{verbatim}

Start with a normal function.

\begin{verbatim}
def show_line(name, score, mood):
    print(f"name={name!r}")
    print(f"score={score!r}")
    print(f"mood={mood!r}")

data = {
    "name": "Homer",
    "score": 42,
    "mood": "hungry",
}

show_line(**data)
\end{verbatim}

Possible output:

\begin{verbatim}
name='Homer'
score=42
mood='hungry'
\end{verbatim}

The call:

\begin{verbatim}
show_line(**data)
\end{verbatim}

means:

\begin{verbatim}
take key-value pairs from data
send them as separate keyword arguments
\end{verbatim}

So this:

\begin{verbatim}
show_line(**data)
\end{verbatim}

is equivalent to:

\begin{verbatim}
show_line(name=data["name"], score=data["score"], mood=data["mood"])
\end{verbatim}

for this specific dictionary.

The dictionary object can be inspected before unpacking.

\begin{verbatim}
data = {
    "name": "Homer",
    "score": 42,
    "mood": "hungry",
}

print(f"type(data): {type(data)}")
print()

selected_names = [
    "keys",
    "values",
    "items",
    "get",
    "__iter__",
    "__getitem__",
]

for name in selected_names:
    print(f"{name:<14} {name in dir(data)}")

print()
print(f"keys:      {list(data.keys())}")
print(f"items:     {list(data.items())}")
print(f"name:      {data['name']!r}")
\end{verbatim}

Possible output:

\begin{verbatim}
type(data): <class 'dict'>

keys           True
values         True
items          True
get            True
__iter__       True
__getitem__    True

keys:      ['name', 'score', 'mood']
items:     [('name', 'Homer'), ('score', 42), ('mood', 'hungry')]
name:      'Homer'
\end{verbatim}

The unpacking operation uses the mapping keys as keyword names.

If a required key is missing, normal argument binding fails.

\begin{verbatim}
def show_line(name, score, mood):
    print(f"name={name!r}")
    print(f"score={score!r}")
    print(f"mood={mood!r}")

data = {
    "name": "Homer",
    "score": 42,
}

try:
    show_line(**data)
except TypeError as err:
    print(f"error type: {type(err)}")
    print(f"message:    {err}")
\end{verbatim}

Possible output:

\begin{verbatim}
error type: <class 'TypeError'>
message:    show_line() missing 1 required positional argument: 'mood'
\end{verbatim}

The mapping did not provide a value for \texttt{mood}.

If the mapping contains an unknown keyword and the function has no \texttt{**kwargs} receiver, binding also fails.

\begin{verbatim}
def show_line(name, score):
    print(f"name={name!r}")
    print(f"score={score!r}")

data = {
    "name": "Homer",
    "score": 42,
    "mood": "hungry",
}

try:
    show_line(**data)
except TypeError as err:
    print(f"error type: {type(err)}")
    print(f"message:    {err}")
\end{verbatim}

Possible output:

\begin{verbatim}
error type: <class 'TypeError'>
message:    show_line() got an unexpected keyword argument 'mood'
\end{verbatim}

If the function has a \texttt{**kwargs} receiver, the unknown keyword can be collected.

\begin{verbatim}
def show_line(name, score, **extra):
    print(f"name={name!r}")
    print(f"score={score!r}")
    print(f"extra={extra}")

data = {
    "name": "Homer",
    "score": 42,
    "mood": "hungry",
}

show_line(**data)
\end{verbatim}

Possible output:

\begin{verbatim}
name='Homer'
score=42
extra={'mood': 'hungry'}
\end{verbatim}

The keyword names from the mapping first try to bind to declared parameters. The leftover keyword names are packed into the variadic keyword dictionary.

The keys used by \texttt{**} must be strings.

\begin{verbatim}
def show_all(**kwargs):
    print(kwargs)

data = {
    "name": "Homer",
    42: "answer",
}

try:
    show_all(**data)
except TypeError as err:
    print(f"error type: {type(err)}")
    print(f"message:    {err}")
\end{verbatim}

Possible output:

\begin{verbatim}
error type: <class 'TypeError'>
message:    keywords must be strings
\end{verbatim}

The double-star call syntax creates keyword arguments, and keyword names must be strings.

Several mappings can be unpacked in one call, but duplicate keyword names are not allowed in a function call.

\begin{verbatim}
def show_all(**kwargs):
    print(kwargs)

first = {"name": "Homer", "score": 42}
second = {"mood": "hungry"}

show_all(**first, **second)
\end{verbatim}

Possible output:

\begin{verbatim}
{'name': 'Homer', 'score': 42, 'mood': 'hungry'}
\end{verbatim}

This works because the keys are distinct.

Duplicate keys fail at call binding time.

\begin{verbatim}
def show_all(**kwargs):
    print(kwargs)

first = {"name": "Homer", "score": 42}
second = {"score": 100, "mood": "focused"}

try:
    show_all(**first, **second)
except TypeError as err:
    print(f"error type: {type(err)}")
    print(f"message:    {err}")
\end{verbatim}

Possible output:

\begin{verbatim}
error type: <class 'TypeError'>
message:    __main__.show_all() got multiple values for keyword argument 'score'
\end{verbatim}

The exact module name may differ. The rule is stable: a function call cannot receive the same keyword name twice.

This is different from dictionary merging, where later values can overwrite earlier values.

\begin{verbatim}
first = {"name": "Homer", "score": 42}
second = {"score": 100, "mood": "focused"}

merged = {**first, **second}

print(f"merged: {merged}")
\end{verbatim}

Possible output:

\begin{verbatim}
merged: {'name': 'Homer', 'score': 100, 'mood': 'focused'}
\end{verbatim}

Inside a dictionary display, duplicate keys are resolved by the resulting dictionary construction. In a function call, duplicate keyword arguments are rejected because a parameter cannot receive two values.

The double-star unpacking operation does not copy the contained objects. It passes references to the same objects.

\begin{verbatim}
def mutate_option(**kwargs):
    items = kwargs["items"]
    items.append("No soup for you!")

quotes = ["D'oh!"]

data = {
    "items": quotes,
}

print("before")
print(f"quotes:     {quotes}")
print(f"id(quotes): {id(quotes)}")
print()

mutate_option(**data)

print("after")
print(f"quotes:     {quotes}")
print(f"id(quotes): {id(quotes)}")
\end{verbatim}

Possible output:

\begin{verbatim}
before
quotes:     ["D'oh!"]
id(quotes): 2292527068032

after
quotes:     ["D'oh!", 'No soup for you!']
id(quotes): 2292527068032
\end{verbatim}

The dictionary \texttt{kwargs} inside the function is newly created for the call, but the value stored under \texttt{"items"} is a reference to the same list object. Mutating that list is visible outside the function.

\subsubsection{Key-Value Parameter Extraction, Mapping Lookups, and Safe Override Patterns}

A \texttt{**kwargs} dictionary is often used for optional named settings. The function can look up keys, provide defaults, reject unknown options, or forward options to another function.

The simplest lookup uses square brackets.

\begin{verbatim}
def show_required(**kwargs):
    name = kwargs["name"]
    score = kwargs["score"]

    print(f"name={name!r}")
    print(f"score={score!r}")

show_required(name="Homer", score=42)
\end{verbatim}

Possible output:

\begin{verbatim}
name='Homer'
score=42
\end{verbatim}

Square-bracket lookup requires the key to exist. If it does not exist, Python raises \texttt{KeyError}.

\begin{verbatim}
def show_required(**kwargs):
    name = kwargs["name"]
    score = kwargs["score"]

    print(f"name={name!r}")
    print(f"score={score!r}")

try:
    show_required(name="Homer")
except KeyError as err:
    print(f"error type: {type(err)}")
    print(f"missing key: {err}")
\end{verbatim}

Possible output:

\begin{verbatim}
error type: <class 'KeyError'>
missing key: 'score'
\end{verbatim}

The safer lookup method for optional keys is \texttt{get()}.

\begin{verbatim}
def show_optional(**kwargs):
    name = kwargs.get("name", "unknown")
    score = kwargs.get("score", 42)
    mood = kwargs.get("mood", "calm")

    print(f"{name:<10} score={score:04d} mood={mood}")

show_optional(name="Homer")
show_optional(name="Lisa", score=100, mood="focused")
\end{verbatim}

Possible output:

\begin{verbatim}
Homer      score=0042 mood=calm
Lisa       score=0100 mood=focused
\end{verbatim}

The method call:

\begin{verbatim}
kwargs.get("score", 42)
\end{verbatim}

means:

\begin{verbatim}
if "score" exists:
    return kwargs["score"]
else:
    return 42
\end{verbatim}

For required-but-manual validation, check membership explicitly.

\begin{verbatim}
def show_checked(**kwargs):
    if "name" not in kwargs:
        raise ValueError("missing required option: name")

    name = kwargs["name"]
    score = kwargs.get("score", 42)

    print(f"{name:<10} score={score:04d}")

show_checked(name="Homer")

try:
    show_checked(score=100)
except ValueError as err:
    print(f"error type: {type(err)}")
    print(f"message:    {err}")
\end{verbatim}

Possible output:

\begin{verbatim}
Homer      score=0042
error type: <class 'ValueError'>
message:    missing required option: name
\end{verbatim}

This kind of check is clearer than allowing a low-level \texttt{KeyError} when the missing key is part of the function's public call contract.

A function can also reject unknown keys.

\begin{verbatim}
def show_limited(**kwargs):
    allowed = {"name", "score", "mood"}

    for key in kwargs:
        if key not in allowed:
            raise ValueError(f"unknown option: {key}")

    name = kwargs.get("name", "unknown")
    score = kwargs.get("score", 42)
    mood = kwargs.get("mood", "calm")

    print(f"{name:<10} score={score:04d} mood={mood}")

show_limited(name="Homer", score=42)

try:
    show_limited(name="Bart", snack="donut")
except ValueError as err:
    print(f"error type: {type(err)}")
    print(f"message:    {err}")
\end{verbatim}

Possible output:

\begin{verbatim}
Homer      score=0042 mood=calm
error type: <class 'ValueError'>
message:    unknown option: snack
\end{verbatim}

The set object used for allowed keys is also an ordinary object.

\begin{verbatim}
allowed = {"name", "score", "mood"}

print(f"type(allowed): {type(allowed)}")
print()

selected_names = [
    "add",
    "remove",
    "discard",
    "union",
    "__contains__",
    "__iter__",
]

for name in selected_names:
    print(f"{name:<14} {name in dir(allowed)}")

print()
print(f"'score' in allowed: { 'score' in allowed }")
print(f"'snack' in allowed: { 'snack' in allowed }")
\end{verbatim}

Possible output:

\begin{verbatim}
type(allowed): <class 'set'>

add            True
remove         True
discard        True
union          True
__contains__   True
__iter__       True

'score' in allowed: True
'snack' in allowed: False
\end{verbatim}

A set is useful here because membership checks are direct and readable.

When extracting options from \texttt{kwargs}, \texttt{pop()} is useful if the function wants to consume recognized options and then check whether anything unknown remains.

\begin{verbatim}
def show_consumed(**kwargs):
    name = kwargs.pop("name", "unknown")
    score = kwargs.pop("score", 42)
    mood = kwargs.pop("mood", "calm")

    if kwargs:
        raise ValueError(f"unknown options: {list(kwargs.keys())}")

    print(f"{name:<10} score={score:04d} mood={mood}")

show_consumed(name="Homer", score=42)

try:
    show_consumed(name="Bart", snack="donut", town="Springfield")
except ValueError as err:
    print(f"error type: {type(err)}")
    print(f"message:    {err}")
\end{verbatim}

Possible output:

\begin{verbatim}
Homer      score=0042 mood=calm
error type: <class 'ValueError'>
message:    unknown options: ['snack', 'town']
\end{verbatim}

This mutates the local \texttt{kwargs} dictionary. That is usually safe because \texttt{kwargs} is a fresh dictionary created for the current function call.

However, if the function receives an ordinary dictionary as a normal parameter, mutating it directly would mutate the caller's dictionary. In that case, make a shallow copy first if the caller's mapping should be preserved.

\begin{verbatim}
def consume_mapping(options):
    local_options = dict(options)

    name = local_options.pop("name", "unknown")
    score = local_options.pop("score", 42)

    print(f"name={name!r}, score={score}")
    print(f"remaining local_options: {local_options}")

settings = {
    "name": "Homer",
    "score": 42,
    "mood": "hungry",
}

print("before")
print(f"settings: {settings}")
print()

consume_mapping(settings)

print()
print("after")
print(f"settings: {settings}")
\end{verbatim}

Possible output:

\begin{verbatim}
before
settings: {'name': 'Homer', 'score': 42, 'mood': 'hungry'}

name='Homer', score=42
remaining local_options: {'mood': 'hungry'}

after
settings: {'name': 'Homer', 'score': 42, 'mood': 'hungry'}
\end{verbatim}

The expression:

\begin{verbatim}
local_options = dict(options)
\end{verbatim}

creates a new dictionary containing the same key-value references. The dictionary container is new. The contained objects are not deeply copied.

A common safe override pattern is to build a new options dictionary before calling another function.

\begin{verbatim}
def show_line(name, score, mood):
    print(f"{name:<10} score={score:04d} mood={mood}")

defaults = {
    "name": "unknown",
    "score": 42,
    "mood": "calm",
}

user_options = {
    "name": "Lisa",
    "score": 100,
}

final_options = {
    **defaults,
    **user_options,
}

print(f"final_options: {final_options}")
show_line(**final_options)
\end{verbatim}

Possible output:

\begin{verbatim}
final_options: {'name': 'Lisa', 'score': 100, 'mood': 'calm'}
Lisa       score=0100 mood=calm
\end{verbatim}

The dictionary construction:

\begin{verbatim}
final_options = {
    **defaults,
    **user_options,
}
\end{verbatim}

means:

\begin{verbatim}
start with defaults
then apply user_options on top
if the same key appears twice, the later value wins
\end{verbatim}

This avoids duplicate keyword errors in the function call because the merging happened before the call. The function receives each keyword only once.

This is unsafe:

\begin{verbatim}
def show_line(name, score, mood):
    print(f"{name:<10} score={score:04d} mood={mood}")

defaults = {
    "name": "unknown",
    "score": 42,
    "mood": "calm",
}

user_options = {
    "name": "Lisa",
    "score": 100,
}

try:
    show_line(**defaults, **user_options)
except TypeError as err:
    print(f"error type: {type(err)}")
    print(f"message:    {err}")
\end{verbatim}

Possible output:

\begin{verbatim}
error type: <class 'TypeError'>
message:    __main__.show_line() got multiple values for keyword argument 'name'
\end{verbatim}

The call attempts to pass \texttt{name} twice and \texttt{score} twice. Python rejects that.

This is safe:

\begin{verbatim}
final_options = {**defaults, **user_options}
show_line(**final_options)
\end{verbatim}

Now the duplicate keys have already been resolved by dictionary construction.

For forwarding options, a function can separate what it consumes from what it forwards.

\begin{verbatim}
def render_text(text, *, uppercase=False, **style):
    if uppercase:
        text = text.upper()

    print(f"text:  {text}")
    print(f"style: {style}")

render_text(
    "No soup for you!",
    uppercase=True,
    font="serif",
    size=14,
    answer=42,
)
\end{verbatim}

Possible output:

\begin{verbatim}
text:  NO SOUP FOR YOU!
style: {'font': 'serif', 'size': 14, 'answer': 42}
\end{verbatim}

The keyword-only parameter \texttt{uppercase} is consumed by the function directly. The remaining keyword arguments are collected into \texttt{style}.

A more explicit forwarding example:

\begin{verbatim}
def final_display(text, **style):
    print(f"text:  {text}")
    print(f"style: {style}")

def display_quote(text, **options):
    uppercase = options.pop("uppercase", False)

    if uppercase:
        text = text.upper()

    final_display(text, **options)

display_quote(
    "It is only a model.",
    uppercase=True,
    font="mono",
    size=12,
)
\end{verbatim}

Possible output:

\begin{verbatim}
text:  IT IS ONLY A MODEL.
style: {'font': 'mono', 'size': 12}
\end{verbatim}

Here, \texttt{display\_quote()} consumes \texttt{uppercase} and forwards the remaining options to \texttt{final\_display()}.

The safe structure is:

\begin{verbatim}
1. receive **options
2. extract options that belong to the current function
3. validate or reject unknown options if necessary
4. forward only the remaining dictionary with **options
\end{verbatim}

The important rule is:

\begin{quote}
\texttt{**kwargs} packs keyword overflow into a dictionary. \texttt{**mapping} at the call site unpacks a mapping into keyword arguments. Dictionary merging can resolve overrides before a call; function-call binding itself rejects duplicate keyword names.
\end{quote}